\label{sec:discussion}

\subsection{Critical evaluation of the design} % REQUIRED HEADING

This subsection interprets the results of Section~\ref{sec:results}, evaluates the design choices in hindsight and identifies what must change before the requirements can be considered proven.
\label{sec:dis-critical}

Two limitations determine every conclusion below: the measurements were taken on a development laptop and not on the ODROID N2+, and the evidence on synthesis quality is partly circular, because the classifier that selects the best of $N$ candidates belongs to the same family as the one that reports the score.
Where a figure is affected by either limitation, this is stated with the figure.

    \subsubsection{Interpretation of results}
    \label{sec:dis-interp}

        \textbf{R2, classification: the strongest part of the system.}
        Text decoding reached 98.61\% character accuracy and 94.21\% word accuracy over 669 real line crops, and writer identification 99.78\% over 1358 real lines from thirteen writers with the ink classifier.
        Both exceed the 95\% target on this evidence, and the lower bound of the character figure is about 97.6\%.
        Two qualifications do not show in the numbers.
        First, the 98.61\% is measured on crops known to be correctly cut, because any page whose crop count disagreed with the label count was excluded; it therefore measures recognition given correct segmentation and says nothing about pages where segmentation fails.
        Second, the five writers scored are also represented in the recogniser's fine-tuning data, so the split favours the system.
        Honestly stated, recognition and identification meet their target on correctly segmented, in-distribution input, and the end-to-end page-to-text accuracy is bounded above by these figures rather than equal to them.
        The gap between character and word accuracy carries information: one wrong character spoils a whole word, so the residual error sits in a minority of words rather than spread thinly.

        \textbf{The writer classifier that is used for real ink is not the one that judges synthesis.}
        The ink classifier reached 99.78\% on real lines and the shape classifier 93.67\%.
        That ordering is the opposite of what the synthesis path needs, because the ink classifier depends on stroke thickness and ink density, which a constant-width plotter cannot reproduce: on the system's own drawn output the ink classifier falls to 28\%, close to chance for thirteen writers, while the shape classifier reaches 89.7\%.
        Both are therefore necessary: one is the right instrument for reading a photograph, the other for judging a drawing.

        \textbf{A training fault, not an intrinsic difficulty, explained the hardest writer pair.}
        Before this work the shape classifier assigned 118 of one writer's 148 real lines to a single other writer, which read as two intrinsically similar hands.
        The measurement shows otherwise: training the same architecture end-to-end, instead of freezing the backbone at the recogniser's features and fitting only the classification head, raised the shape classifier from 83.4\% to 93.7\% overall and that writer from 18.9\% to 67.6\%.
        A recogniser's features are trained to be \emph{invariant} between writers, so a classifier restricted to them can separate only writers whose recognition features already differ.
        The residual is still 48 of 148 lines sent to the same neighbour, so a smaller genuine effect probably remains on top of the fault; this is unresolved and listed as an unsolved problem.
        The lesson is that reusing a recogniser's backbone saves training data, but it must be fine-tuned rather than frozen whenever the new task is to detect the variation the backbone was trained to ignore.

        \textbf{R1, representation fidelity: close, and the shortfall is in the right place to be fixable.}
        Synthesised text was read at 82.2\% character accuracy over thirteen writers, against the 85\% target, and at 79.2\% in a wider coverage run.
        The matching comparison is more informative: on identical sentences the writers' genuine handwriting was read at 89.0\% and the synthesis at 81.6\%.
        The synthesis loses 7.4 percentage points of character accuracy but 23.8 points of word accuracy, and that asymmetry localises the defect.
        Individual letterforms survive the transfer, as expected because the synthesiser draws the writer's own traced glyph paths; the spacing structure does not, so words are drawn with boundaries the recogniser reads differently.
        The measured feature deviations agree: synthesised lines have ascenders about 0.4 to 0.6 x-height units shorter than the writer's real lines, word spaces about 0.5 units narrower, and connectedness lower.
        The defect is thus a few measurable, correctable quantities rather than a failure of the method, which is why R1 is partially met rather than not met.
        A recogniser checkpoint that took no part in the search read the same lines 3.7 points below the production checkpoint, so the selection raises legibility only slightly and the result is not an artefact of the search.

        \textbf{The stylistic-correlation target is not verified, and writer identification is a weak proxy for it.}
        The proposal asks for 85\% stylistic correlation; what was measured is writer-identification accuracy, a different question, because a synthesis could be identifiable as the right writer while still differing visibly from that writer's hand.
        Worse, the writer-identification figures on synthesised text (80.1\% in the official run, 89.7\% in the coverage run) come from a classifier that also selects the candidates, so a best-of-$N$ search necessarily raises the score it is judged by, and only the figures computed outside the search can be trusted.
        The nine-feature geometric correlation gave 0.714 against the 0.85 target, so the style term of R1 is not met on that measure, the shortfall being in descender extent, the ascender-to-descender ratio and word spacing rather than in slant or ascender extent.
        The overlay test gave an AUC~\cite{fawcett2006} of about 0.71, against 0.5 for an output carrying no writer information, which indicates real but limited fidelity.
        A human panel or an independent classifier is required before the 85\% stylistic target can be claimed, and this was not carried out.

        \textbf{The synthesis is recognised as the correct writer much better than the system draws that writer.}
        The same shape classifier identifies 93.7\% of real lines but only 89.7\% of drawn lines in the coverage run, and only 67.6\% of one writer's real lines while identifying 100\% of that writer's drawn lines.
        A judge that recognises the machine's output more reliably than the human's hand is measuring agreement with the synthesiser rather than fidelity to the author, which is the most important caveat on every writer-identification figure measured on synthesised text.

        \textbf{R3, reproduction time: the failure is the search, not the drawing.}
        A four-character word at $N=6$ took 64.3~s, or 16.1~s per character, four times the limit.
        The coverage run confirms that the cost scales with the number of candidates, because every candidate is rendered, read by the recogniser and scored by the writer classifier.
        The physical drawing time was never measured, so the requirement cannot be decomposed, and it would be wrong to conclude that the gantry is slow: the evidence points at the selection procedure, which embeds two network passes in its inner loop and pays for the quality it delivers.
        The proposal does not state whether the synthesis time counts towards the limit; under the reading that only the physical writing counts, the generator's model value of 2.7 to 4.0~s per character sits at the 4~s limit, so the requirement is recorded as not met on the conservative reading that the total counts.

        \textbf{R4, classification time: recognition passes, segmentation fails badly.}
        Recognition ran at 8.86~ms per character against the 25~ms allowance, a margin of about 2.8 times, so the neural component is not the problem.
        Page segmentation took 35 to 57~s against the roughly 2~s the equation allows for capture and conditioning, a shortfall of more than an order of magnitude.
        Since $x$ appears only in the recognition term, a longer page does not buy proportionally more time, so the segmenter's per-page cost cannot be amortised.
        The profile shows the cause: line grouping and the chain search dominate, at about 70\% of the time, and are written in pure Python.
        Moving two inner routines to C reduced the total by 25\% to 32\%, but the grouping still dominates: accelerating the wrong part of a program cannot fix it.

        \textbf{R5, positional accuracy, and the field conditions: no conclusion is possible.}
        The derived step resolution of 12.5~$\mu$m shows only that the command quantisation is 40 times finer than the tolerance, so the error must come from the mechanics and from open-loop effects.
        No drawn mark was scanned, because the direction-pin fault prevented controlled motion in both directions, so the requirement is unassessed rather than failed.
        The same holds for the two field conditions, neither of which was measured; the illumination figure is the more consequential omission, because the conditioning stage was designed to tolerate uneven lighting and that claim is untested.

        \textbf{Requirements met.}
        The storage half of R6 is met: thirteen styles are stored as files and the live server lists all of them, against a minimum of ten, returning the list in 70~ms. The reproduction half is not: the proposal's demonstration asks for a writer to be chosen at random and the text written on paper in that style, and that was never carried out.
        R2 is met for writer identification on real handwriting and for recognition on correctly segmented in-distribution crops, but not demonstrated for the complete page-to-text chain.
        R1 is close but not met, R3 and R4 are not met, and R5 and both field conditions are unassessed.


    \subsubsection{Critical evaluation}
    \label{sec:dis-crit}

        \textbf{Line-level reading instead of character separation.}
        The design replaced the proposal's character separation by separation into lines and a sequence recogniser.
        This was right for connected handwriting, where character boundaries do not exist, and it makes the accuracy independent of a separation error; the price is that a whole line must be read by one network, at a cost per character higher than the 25~ms the proposal assumed.

        \textbf{Best-of-$N$ search with the judges inside the loop.}
        The search with a recogniser and a writer classifier as judges was good for quality, because it prevents the synthesis from drawing an illegible or wrong-style sample.
        In hindsight it was poor for verification, because it removed the independence of the evaluation, and for time, because each of the $N$ candidates costs a rendering and two network passes.
        The qualification protocols address the first, and a lower $N$ or a cheaper first-stage filter the second.

        \textbf{Pure Python for the segmentation.}
        The proposal required that the processing run on the embedded platform and be written without library functions for the algorithm itself.
        Writing the image processing in Python made the design transparent and correct, but the inner loops of the line grouping and labelling are too slow.
        The C routines reduce the labelling and the box sums by 84\% to 88\%, but the end-to-end effect (25\% to 32\%) is small because the grouping dominates.
        By Amdahl's law, if a fraction $p$ of the time is accelerated by a factor $s$, the total speed-up is $1/((1-p)+p/s)$, so with $p=0.30$ (the share outside the grouping, taking the grouping as 70\%) it cannot exceed $3.3$ even if the accelerated part took no time.
        The acceleration was not deployed, because the benefit was not decisive.

        \textbf{Open-loop gantry with step counting.}
        Using the stepper motors without encoders is a standard, low-cost choice; its accuracy depends on missed steps, and no sensor reveals them.
        Running the driver without acceleration ramps (the generator models an acceleration of 400~mm/s$^2$, the driver does not execute it) makes the real motion differ from the model and a lost step more likely at the start of a stroke at higher feed.


    \subsubsection{Unsolved problems}
    \label{sec:dis-unsolved}

        The following problems were not solved.

        \begin{enumerate}
          \item \textbf{X-axis direction fault.} During bring-up the X-axis calibration reached the first end-stop, but the direction pin (physical pin 7) was not observed to change level when the direction was reversed.
          The fault was unresolved when the project closed, and no cause is claimed, because the evidence does not support one.
          Qualification tests 3 (writing time), 5 and 6 (physical demonstration) depend on its resolution.
          \item \textbf{No measurements on the target hardware.} No timing on the ODROID N2+ and no physical accuracy measurement exists; the figures for R3 and R4 are from a laptop.
          \item \textbf{The undertaking to process on the embedded platform was not demonstrated.} The approved proposal states, as a ticked condition, that ``all processing will be done on an embedded platform, and not on a PC platform''.
          The deployed programs do run on the ODROID N2+, and the inference code needs nothing the board cannot provide, but every measurement in Section~\ref{sec:results} was taken on a laptop, and training was done off the board entirely.
          The undertaking is therefore met in the design and unproven in the evidence: the most significant gap between what the proposal promised and what this report can show, and a measurement gap rather than a design gap, since the same programs need only be run on the board and timed.
          \item \textbf{R4 not met, and the shortfall is entirely in the segmentation.} The segmentation time exceeds the two seconds the specification allows for capture and conditioning by more than an order of magnitude. The recognition time does not: at 8.86~ms per character it sits a factor of about 2.8 \emph{inside} the 25~ms allowance. The requirement fails on the conditioning stage alone.
          \item \textbf{Optimistic validation of the personal writers.} A line-level split inside pages was used, and the checkpoint was selected on the reporting data; no separate test set exists.
          \item \textbf{Circular synthesis evaluation.} No independent judge or human panel was used.
          \item \textbf{Architecture sweep.} A sweep over widths, depths and learning rates was programmed but no result was kept, so the network layers are those of the published design.
          \item \textbf{Implementation issues of the writing path.} (i) The stroke ordering sorts text rows by ascending $Y$, and with $Y$ upward this writes the bottom line first, while the documentation says reading order. (ii) The acceleration the generator models is not executed by the driver. (iii) The generator's work area (200~mm) is larger than the driver's travel between its limits (194~mm in X). (iv) The generator issues the fast feed once and the driver has a single modal feed register, so pen-up travel runs at the pen-down feed (900~mm/min instead of 2400~mm/min), lengthening the job.
          \item \textbf{Missing information.} The hardware documentation (CAD files, photographs, wiring diagram, camera data) and the inconsistent values of the number of tries (14 in the tuning notes, 20 as the server default, 30 mentioned in the interface) are \mustdo{reconcile the values and supply the hardware documentation}.
          \item \textbf{Training mode.} The proposal planned training on the single-board computer; training is done offline, and the interface has no training function.
          \item \textbf{Display of the classification time} (specified in R4) is not implemented.
          \item \textbf{The CTC loss is not safe against an empty transcription.} The gradient check of Section~\ref{sec:textrec-numpynn} found that for a label sequence of length zero the recursion builds two of its shifted arrays at the wrong length and raises a shape error instead of returning a loss.
          Every sequence of length one or more passes, including the repeated-character case that exercises the forbidden skip transition, so the derivative is correct and this is an indexing fault at the boundary.
          It was left unfixed because no line in the data set has an empty transcription, but a future data set would expose it; the same check showed that samples for which no alignment exists must be filtered out before training rather than left to the loss to reject.
        \end{enumerate}

        \textbf{What must change to meet R4.}
        Three measures are required together.
        First, the line grouping (about 70\% of the segmentation time) must be redesigned, for example by replacing the exhaustive search over component pairs by neighbourhood queries on a spatial index, if the profile on the ODROID N2+ confirms that the search is the quadratic part; the target is a reduction by at least a factor of 20, because $32~\mathrm{s}$ (70\% of 46~s) must fall well below 2~s.
        Second, the image must be reduced in resolution before segmentation: halving both dimensions of the $2400\times1800$ photograph reduces the pixel count by 4, and, as the filters and labelling scale with the pixel count, the share outside the grouping (about 14~s) should fall to about 3.5~s, and to about 0.4 to 0.6~s with the C routines, which are 84\% to 88\% faster (a factor of 6 to 8.3).
        These are projections from the measured ratios, not measurements, and they ignore that the ODROID N2+ is slower than the laptop.
        Third, the recogniser must become about 3 times cheaper per character, for example by a smaller network, by quantising the weights to 8 bit integers or by reducing the output sequence length; the cost per character is roughly proportional to the multiply-accumulate operations per line, and each measure must be validated against the accuracy loss, since R2 is already under its target.

        \textbf{What must change for the other requirements.}
        For R2 (text): more personal training pages, an evaluation on held-out pages, and a language-model or lexicon decoder.
        For R1: an independent evaluation and the human panel, and re-tuning with an independent judge if the independent figure is lower.
        For R3: measurement on the ODROID N2+, correction of the feed register, execution of acceleration or a lower pen dwell, and a lower $N$.
        For R5: resolution of the direction-pin fault, completion of the calibration and execution of qualification test~5.

    \subsubsection{Strong points of the design}
    \label{sec:dis-strong}

        Each strong point follows from a design decision.
        \begin{itemize}
          \item Reading whole lines removes the dependence on a character separation that cannot work for connected handwriting, and the recogniser reached 91.82\% on an independent general test set.
          \item The conditioning stage (flat-field correction, local thresholding, skew correction) allows a photograph with uneven lighting to be used, and the live check on a personal photograph read the correct writer at 98.6\% confidence with 28 lines detected and an accurate transcription.
          \item Writers are compact profile files, so a new writer is added by a file, and thirteen writers are stored permanently, satisfying the storage half of R6.
          \item The two judges in the search (text and writer) give a measurable improvement over a single sample, and the shape classifier solves the problem that an ink-based classifier cannot judge a plotter's constant line width (about 20\% on synthesised output).
          \item The numerical resolution of the gantry (12.5~$\mu$m) is 40 times finer than the tolerance of R5, and the absolute step bookkeeping prevents accumulating rounding errors.
          \item Safety in software: end-stop confirmation (three consecutive samples), software limits and automatic retraction after a trip.
          \item The algorithms are implemented from first principles and documented well enough to be analysed, and the known limitations (circularity, optimistic splits) are stated.
        \end{itemize}

    \subsubsection{Expected failure conditions}
    \label{sec:dis-failure}

        The system is expected to fail or degrade strongly under the following conditions.
        \begin{itemize}
          \item \textbf{Lighting.} Illuminance below 400~lux or strongly directional light with shadows violates FC1; the thresholding then merges ink with shadows or loses faint strokes, and recognition accuracy falls.
          \item \textbf{Page skew and perspective.} Large skew angles or a page that is not flat exceed the range of the skew correction.
          \item \textbf{Non-text regions.} Diagrams, tables and printed text can be segmented as text lines; the text versus non-text scoring reduces but does not exclude this, and the recogniser then produces meaningless characters.
          \item \textbf{Unseen writers.} The writer classifier is a thirteen-class closed-set classifier, so a writer who is not among the thirteen is assigned to one of them, because the softmax cannot express ``none of the above'': the output is wrong without a warning.
          \item \textbf{Styles outside the glyph library.} A writer whose letter shapes differ greatly from the stored paths (unusual letters, capitals, digits or symbols that were scarce in the training pages) is reproduced with the nearest stored shape; characters never harvested cannot be written in that style.
          \item \textbf{Text with accents, mathematical notation or rare symbols}, which occur in the personal pages, are outside the benchmark vocabulary.
          \item \textbf{Motor stalls and lost steps.} The open-loop gantry has no feedback, so a stall shifts every later mark without any signal.
          \item \textbf{Limit switch failure.} A switch that does not close lets the carriage run against the mechanical stop; the software protection depends on the switches, and a hardware emergency stop that cuts the power does not exist.
          \item \textbf{Long sentences.} The synthesis time grows with the number of characters and tries, so a long text at a high quality setting would take minutes on the board and may time out the web request.
          \item \textbf{Network failure} between the browser and the server prevents control of the system; a running job continues on the board.
        \end{itemize}
        The system does not fail gracefully in all cases: the classifier returns a confident answer for an unseen writer, and the gantry cannot detect lost steps.

\subsection{Considerations in the design} % REQUIRED HEADING
\label{sec:dis-considerations}

This subsection addresses the ergonomics, safety, environmental, social, legal and ethical aspects of the design.

    \subsubsection{Ergonomics}

        The user works through a web interface of four pages (Read, Write, Model and Style) in a browser on any laptop on the network, so nothing is installed.
        The Read page captures the page with a live preview and shows the recognised text, the writer and the confidence; the Write page keeps the writing grid locked until the gantry is calibrated and shows a G-code preview before sending.
        The shortcomings are that the live progress of a running job is not displayed, the classification time is not shown, and the interface has no training function.
        The rig is set up by hand, with the page, the pen and the paper placed by hand; its documentation is incomplete.

    \subsubsection{Health and safety}

        The electronics operate at low voltage (a single-board computer, stepper-driver modules and small motors), so the risk of electric shock is low; \mustdo{supply voltage and the mains insulation of the power supply}.
        The mechanical hazards are the moving carriage, belts and pulleys, which are pinch points while the gantry moves; the pen tip is not hazardous at the writing force.
        The protective measures are the end-stop switches with confirmation, the software limits, the automatic retraction and calibration at a low speed (7.8~mm/s).
        There is no hardware emergency stop that removes the motor power, and the server has no abort function for a running job; a job can be stopped only by ending the program or removing the power, so fitting a mains-independent emergency stop is recommended for use by others.
        Users must keep their hands away from the carriage while a job runs.

    \subsubsection{Environmental impact}

        The ODROID N2+ is a low-power board and the stepper motors draw power only while the gantry works, so the energy use is small beside a desktop computer.
        The system uses paper; test pages and failed jobs add waste, limited by the preview before writing and by the option to write on used paper.
        The mechanical parts printed in plastic or bought as standard parts are \mustdo{which parts were 3D-printed and with which material}; they can be reused, and plastic and electronic waste must go to a recycling point, not into domestic waste.
        The pen ink is a small consumable, no hazardous material beyond the electronics is used, and the only emissions are the motor noise and the pulses of the drive.

    \subsubsection{Social and legal impact}

        The social benefit is that the system can write in a person's own handwriting for those who cannot write neatly or at all, for example after an injury, and offers the personal touch the proposal identified.
        There is also a risk of misuse: reproducing a person's handwriting can serve forgery or impersonation, for instance on signed documents or letters.
        The plotter pen's constant line width reduces but does not remove the risk, because it differs from natural writing in stroke width and pressure (the ink-based classifier identified only about 20\% of the synthesised output, so these ink properties are not reproduced), yet a document in a plausible style can still mislead a reader.
        Handwriting samples are personal information and a writer profile is similar to biometric data, since it permits identification of the writer; in South Africa the Protection of Personal Information Act (POPIA) applies to such data, requiring the consent of the writers, the minimum storage and the protection of the profiles against access by others.
        The samples of the five personal writers require documented consent, \mustdo{confirm that consent was given and in what form}, and the eight database writers contributed to a public research database whose terms of use apply.
        The server has no user accounts (only an optional shared key), accepts requests from any origin and can start physical motion for anyone who can reach its port, so it must be operated only on a trusted local network.
        Beyond the general electrical safety rules no other legislation was identified for a research prototype of this type; \mustdo{any applicable regulation or university rule}.

    \subsubsection{Ethics clearance}

        The project used handwriting samples of five named persons and of a public database, so it involves human participants in a limited way.
        The status of ethics clearance is \mustdo{ethics clearance number, or a statement that no clearance was required}.
        If the human panel of qualification test~1 is carried out, the raters give informed consent and their responses are stored without personal identifiers, and clearance must be confirmed before the panel is run.

